\SourcePage{302}{305}
\section{Second quantization: Fermi statistics}

The essential basis of the method of second quantization remains unchanged
for systems of identical fermions. The specific formulas for matrix elements
of quantities and for the operators $\hat a_i$, however, naturally change.

\SourcePage{303}{306}
The wave function $\psi_{N_1N_2\ldots}$ now has the form (61.5). Its
antisymmetry first raises the question of choosing its sign. No such question
arose with Bose statistics: because the wave function was symmetric, a sign
once chosen was retained under every permutation of the particles.

To fix the sign of (61.5), adopt the following convention. Number all the
states $\psi_i$ once and for all in sequence. The rows of the determinant
(61.5) are then always filled so that
\begin{equation}
 p_1<p_2<p_3<\cdots<p_N,
 \tag{65.1}
\end{equation}
while its columns contain functions of the different variables in the order
$\xi_1,\xi_2,\ldots,\xi_N$. No two of $p_1,p_2,\ldots$ can be equal, since
the determinant would otherwise vanish. Equivalently, every occupation number
$N_i$ can have only the value 0 or 1.

Consider again an operator of the form (64.1),
$\hat F^{(1)}=\sum_a\hat f_a^{(1)}$. For the same reasons as in \S~64, its
matrix elements can be nonzero only for transitions in which all occupation
numbers remain unchanged, or in which one of them, $N_i$, decreases by one
(from unity to zero) while another, $N_k$, increases by one (from zero to
unity). For $i<k$ one readily finds
\begin{equation}
 \langle1_i,0_k|\hat F^{(1)}|0_i,1_k\rangle
 =f_{ik}^{(1)}(-1)^{\sum(i+1,k-1)}.
 \tag{65.2}
\end{equation}
Here $0_i,1_i$ denote $N_i=0,N_i=1$, and $\sum(k,l)$ denotes the sum of the
occupation numbers of all states from $k$ through $l$.\footnote{For $i>k$,
the exponent in (65.2) must be $\sum(k+1,i-1)$. For $i=k\pm1$, these sums are
to be replaced by zero.}
\[
 \sum(k,l)=\sum_{n=k}^lN_n.
\]
The diagonal elements are still given by (64.4):
\begin{equation}
 \overline{F^{(1)}}=\sum_i f_{ii}^{(1)}N_i.
 \tag{65.3}
\end{equation}

For $\hat F^{(1)}$ to be representable in the form (64.13), the annihilation
operators $\hat a_i$ must be defined as matrices

\SourcePage{304}{307}
with elements
\begin{equation}
 \langle0_i|\hat a_i|1_i\rangle
 =\langle1_i|\hat a_i^+|0_i\rangle=(-1)^{\sum(1,i-1)}.
 \tag{65.4}
\end{equation}
Multiplication of these matrices gives, for $k>i$,
\[
 \begin{aligned}
 \langle1_i,0_k|\hat a_i^+\hat a_k|0_i,1_k\rangle
 &=\langle1_i,0_k|\hat a_i^+|0_i,0_k\rangle
   \langle0_i,0_k|\hat a_k|0_i,1_k\rangle\\
 &=(-1)^{\sum(1,i-1)}
   (-1)^{\sum(1,i-1)+\sum(i+1,k-1)},
 \end{aligned}
\]
or
\begin{equation}
 \langle1_i,0_k|\hat a_i^+\hat a_k|0_i,1_k\rangle
 =(-1)^{\sum(i+1,k-1)}.
 \tag{65.5}
\end{equation}
For $i=k$, the matrix $\hat a_i^+\hat a_i$ is diagonal, with element unity
when $N_i=1$ and zero when $N_i=0$. Thus
\begin{equation}
 \hat a_i^+\hat a_i=N_i.
 \tag{65.6}
\end{equation}
Substitution of these expressions in (64.13) indeed gives (65.2), (65.3).

Multiplying $\hat a_i^+$ and $\hat a_k$ in the opposite order gives
\[
 \begin{aligned}
 \langle1_i,0_k|\hat a_k\hat a_i^+|0_i,1_k\rangle
 &=\langle1_i,0_k|\hat a_k|1_i,1_k\rangle
   \langle1_i,1_k|\hat a_i^+|0_i,1_k\rangle\\
 &=(-1)^{\sum(1,i-1)+\sum(i+1,k-1)+\sum(1,i-1)+1},
 \end{aligned}
\]
or
\begin{equation}
 \langle1_i,0_k|\hat a_k\hat a_i^+|0_i,1_k\rangle
 =-(-1)^{\sum(i+1,k-1)}.
 \tag{65.7}
\end{equation}
Comparison of (65.7) with (65.5) shows that these quantities have opposite
signs, and hence
\[
 \hat a_i^+\hat a_k+\hat a_k\hat a_i^+=0,\qquad i\ne k.
\]
For the diagonal matrix $\hat a_i\hat a_i^+$ we find
\begin{equation}
 \hat a_i\hat a_i^+=1-N_i.
 \tag{65.8}
\end{equation}
Adding (65.6) gives
\[
 \hat a_i\hat a_i^++\hat a_i^+\hat a_i=1.
\]
Both relations can be written together as
\begin{equation}
 \hat a_i\hat a_k^++\hat a_k^+\hat a_i=\delta_{ik}.
 \tag{65.9}
\end{equation}
An analogous calculation for products of $\hat a_i$ and $\hat a_k$ yields
\begin{equation}
 \hat a_i\hat a_k+\hat a_k\hat a_i=0
 \tag{65.10}
\end{equation}
(in particular, $\hat a_i\hat a_i=0$).

\SourcePage{305}{308}
Thus operators $\hat a_i$ and $\hat a_k$ (or $\hat a_k^+$) with $i\ne k$
anticommute, whereas under Bose statistics they commuted. This distinction is
natural. In the Bose case, $\hat a_i$ and $\hat a_k$ were completely
independent: each $\hat a_i$ acted only on the single variable $N_i$, and its
action did not depend on the other occupation numbers. Under Fermi statistics,
by contrast, the action of $\hat a_i$ depends both on $N_i$ itself and, as
(65.4) shows, on the occupation numbers of all preceding states. The actions
of different $\hat a_i,\hat a_k$ therefore cannot be regarded as independent.

Once the properties of $\hat a_i,\hat a_i^+$ have been defined, all the other
formulas (64.13)--(64.18) remain fully valid. So do (64.23)--(64.25), which
express operators of physical quantities through the field operators defined
by (64.20). The commutation rules (64.21), (64.22), however, are replaced by
the anticommutation relations
\begin{equation}
 \hat\psi^+(\xi')\hat\psi(\xi)+\hat\psi(\xi)\hat\psi^+(\xi')
 =\delta(\xi-\xi'),
 \tag{65.11}
\end{equation}
\begin{equation}
 \hat\psi(\xi')\hat\psi(\xi)+\hat\psi(\xi)\hat\psi(\xi')=0.
 \tag{65.12}
\end{equation}

If a system contains different kinds of particles, separate second-quantized
operators must be introduced for each kind, as mentioned at the end of the
preceding section. Operators belonging to bosons commute with those belonging
to fermions. Within nonrelativistic theory, operators belonging to different
fermions may formally be taken either to commute or to anticommute; the method
of second quantization gives the same results under either convention.

With a view to later use in relativistic theory, which permits transformations
between different particles, the creation and annihilation operators of
different fermions must nevertheless be taken to anticommute. This becomes
clear if two different internal states of the same composite particle are
regarded as different particles.
